\documentclass{article}
\usepackage{amsmath}

\title{Greedy Basics}
\author{Kanishk Goel}
\date{2 February 2026}
\parindent 0px
\begin{document}
	\maketitle
	\section{Theme}
	Construct a solution iteratively, via a sequence of myopic decisions, and hope that everything works out in the end.
	\section{Scheduling Problem}
	\begin{center}
		\begin{large}
			\textbf{Completion Times} \\
		\end{large}
	\end{center} The \textit{complete time} $C_j(\sigma)$ of a job $j$ in a schedule $\sigma$ is the sum of the lengths of the jobs preceding $j$ in $\sigma$, plus the length of $j$ itself 
	
	\begin{center}
		\begin{large}
			\textbf{Objective Function:} \\
		\end{large}
	\end{center}
	\[\min\limits_{\sigma} \sum_{j=1}^{n} w_j C_j(\sigma)\]
	i.e  sum of weighted completion times. \\
	
	\begin{center}
		\begin{large}
			\textbf{Problem: Minimizing the Sum of Weighted Completion Times} \\
		\end{large}
	\end{center}
	
	\textbf{Input:}  A set of $n$ jobs with positive lengths $l_1,l_2,\ldots,l_n$ and positive weights $w_1,w_2,\ldots,w_n$ \\
	
	\textbf{Output:} A job sequence that minimizes the sum of weighted completion times. \\
	
	\begin{center}
		\begin{large}
			\textbf{Idea} \\
		\end{large}
	\end{center}
	
	\begin{itemize}
		\item If the weights are same and lengths are different, then longer length should get longer completion times (ascending order).
		\item If the lengths are same then, bigger weights should get shorter completion time (descending order).
	\end{itemize}
	
	Therefore, the scores possible for a job are: 
	\begin{itemize}
		\item \textbf{GreedyRatio} -  $j: \frac{w_j}{l_j}$ 
		\item \textbf{GreedyDiff} - $j: w_j - l_j$
	\end{itemize}
	Higher the score, lower the scheduling postion\ldots \\
	
	Let Job $1$ have $l_1 = 5$ and $w_1 = 3$, and Job $2$ have $l_2 = 2$ and $w_2 = 1$. If you place $1$ first and $2$ later, then completion times are $3*5 + 1*(5+2) = 22$ and vice versa it would be $1*2 + 3*(5+2) = 23$. Therefore, $1$ should be placed before $2$. \\
	
	Greedy Differences are $-2$ and $-1$ respectively, which shows GreedyDiff algorithm is wrong, while the Greedy Ratios are $\frac{3}{5} \land \frac{1}{2}$ which schedules the jobs correctly.
	
	\begin{center}
		\begin{large}
			\textbf{Proof of Correctness} \\
		\end{large}
	\end{center}
	
	 Core plan - use \textit{exchange arguments}.  \\
	
	The key idea is to prove that every feasible solution can be improved by modifying it to look more like the output of the greedy algorithm.
	
	\begin{center}
		\begin{small}
			\textbf{Two Assumptions} \\
		\end{small}
	\end{center}
	
	\begin{enumerate}
		\item The jobs are indexed in nonincreasing order of Greedy Ratios
		\item There are no ties between ratios i.e $\frac{w_i}{l_i} \neq \frac{w_j}{l_j}$
	\end{enumerate}
	
	Now the greedy schedule, $\sigma$ is simply the order of jobs $1,2\ldots n$. Therefore $\sigma*$ differs from $\sigma$ because it has some pair of $(i,j) $such that $J_i > J_j$ while $i < j$, where $J_z$ denotes the $z^{th}$ job value as per the greedy schedule.  \\
	
	\textbf{Lemma:} every non greedy schedule has an inversion. \\
	
	Let there be a pair of $(i,j)$ such that $i$ and $j$ are consecutive in a schedule, $\sigma*$ and $C_b$ be completion time of stuff before, we isolate and see for differing portion of $C_{\sigma*}$ and $C_{\sigma}$:
	
	\begin{align}
		\frac{w_i}{l_i} &< \frac{w_j}{l_j} \notag \\
		w_i \cdot l_j &< w_j \cdot l_i \\
		C_{\sigma*}  &= w_i \cdot (C_b + l_i) + w_j \cdot (C_b + l_i + l_j) \\
		C_{\sigma} &= w_j \cdot (C_b + l_j) + w_i \cdot (C_b + l_i + l_j)
	\end{align}
	
	By $(2)$ and $(3)$, we see:
	\begin{align*}
		C_{\sigma*} - C_{\sigma}  &= w_j \cdot l_i - w_i \cdot l_j 
	\end{align*}
	
	Since, RHS is positive, this means $C_{\sigma} < C_{\sigma*}$. Hence, contradiction has been established.
	
	
\end{document}